% Blueprint note: Fixed upstream source versions: \cite{KIPSourcebbac1fc05843}. Interface scope notes are retained as native hgraph node comments; historical migration records are local archives.

% Mathematical entry points on the fixed sphere and synthetic comparison system.
% Paper-labelled conclusions are bound in ../paper/Main.tex instead of duplicated here.
\section{The fixed sphere and its obstruction classes}
Let $E_r^{s,t}$ be the classical Adams spectral sequence of the sphere from Chapter 2, with stem $t-s$. Write $S^{0,0}$ for the synthetic sphere, with its quotient maps, comparison maps and convergence filtration, and set
\[
 W=x_{126,8,4}+x_{126,8},\qquad U=h_0^2x_{124,8},\qquad T=h_1h_4x_{109,12}.
\]
Here the classes are as in Chapter 4. A class $x$ detects a homotopy class $a$ when the image of $a$ in the corresponding associated graded quotient equals the image of $x$.

\begin{definition}[The only possible last differential]
\label{sem:Main:last-differential}
\level{coarse}
\lean{KIP126.Main.H6SquareSupportsD12}
From \cite{LWX126}, Proposition 7.8, statement (2).\par
Let $D$ be the assertion that there are compatible twelfth-page representatives of $h_6^2$ and $T$ with nonzero target and differential
\[
 d_{12}(h_6^2)=T,\qquad (2,128)\longrightarrow(14,139).
\]
\end{definition}



\begin{definition}[Vanishing of the sixth differential]
\label{sem:Main:sixth-differential}
\level{coarse}
\lean{KIP126.Main.Stem126D6Vanishes}
From \cite{LWX126}, Proposition 7.8, statement (3).\par
Let $C_3$ be the assertion that $W$ has a sixth-page representative satisfying
\[
 d_6(W)=0,\qquad (8,134)\longrightarrow(14,139).
\]
\end{definition}

\begin{definition}[Admissible theta and U representatives]
\label{sem:Main:admissible-representatives}
\level{coarse}
\lean{KIP126.Main.ThetaChoice,KIP126.Main.UChoice}
From \cite{LWX126}, Notation 7.2 and Proposition 7.8.\par
An element $\theta\in\pi_{62,64}S^{0,0}$ is admissible if it is nontrivially detected by $h_5^2$. An element $u\in\pi_{124,134}S^{0,0}$ is admissible if it is nontrivially detected by $U$. In both cases, detection uses the comparison and convergence maps of the synthetic sphere. Representatives may differ by a class of higher filtration.
\end{definition}

\begin{definition}[The square-detection condition]
\label{sem:Main:square-detection}
\level{coarse}
\lean{KIP126.Main.Theta5SquareNonzeroDetectionAt,KIP126.Main.Theta5SquareNonzeroDetection}
From \cite{LWX126}, Proposition 7.8, statement (4).\par
For an admissible $\theta$, let $C_4(\theta)$ say that $\theta^2\in\pi_{124,128}S^{0,0}$ is nontrivially detected by $\lambda^6 U$, in tridegree $(10,134,128)$. The condition $C_4$ asserts that $C_4(\theta)$ holds for some admissible $\theta$.
\end{definition}

\begin{definition}[The eta-extension detection condition]
\label{sem:Main:eta-detection}
\level{coarse}
\lean{KIP126.Main.EtaExtensionDetectionAt,KIP126.Main.EtaExtensionDetection}
From \cite{LWX126}, Proposition 7.8, statement (5).\par
For an admissible $u$, let $C_5(u)$ say that $\lambda^3\eta u\in\pi_{125,133}S^{0,0}$ is detected by $\lambda^6 T$, in tridegree $(14,139,133)$. The condition $C_5$ asserts that $C_5(u)$ holds for some admissible $u$.
% Blueprint note: This predicate imposes associated graded detection, rather than an exact representative equation or an additional nonzero-target hypothesis.
\end{definition}

\begin{definition}[The nonzero permanence endpoint]
\label{sem:Main:nonzero-permanence}
\level{coarse}
\lean{KIP126.Main.H6SquarePermanent}
From \cite{LWX126}, Theorem 7.1.\par
Let $P$ assert nonzero survival of $h_6h_6\in E_2^{2,128}$: there exists a representative in $Z_\infty$ mapping to $h_6h_6$ whose image in $E_\infty$ is nonzero.
% Blueprint note: The nonzero-survival formulation retains a common permanent representative and its nonzero limiting image.
\end{definition}

\section{The second-line classification target}
\begin{definition}[The exact second-line classification proposition]
\label{sem:Main:two-line-classification-goal}
\level{coarse}
\lean{KIP126.Main.AdamsTwoLineListed,KIP126.Main.AdamsTwoLineClassification}
From \cite{LWX126}, Corollary 1.5.\par
For $x\in E_2^{2,t}$, let $L(t,x)$ mean that $x$ equals one of
\[
 h_0h_2,\quad h_0h_3,\quad h_2h_4,\quad h_1h_j\ (j\geq3),\quad h_j^2\ (0\leq j\leq6).
\]
Each equality includes its bidegree equation and transport in the actual sphere $E_2$ module. Define the proposition $C$ by
\[
 C\;:\!\Longleftrightarrow\;\forall t\in\mathbb Z\ \forall x\in E_2^{2,t},\quad
 x\text{ survives nontrivially to }E_\infty\ \Longleftrightarrow\ L(t,x).
\]
This entry defines the target; it does not prove $C$.
% Blueprint note: Only proposition-valued definitions are bound here. This entry specifies only the proposition. Its value is separately proved by adams_two_line_classification, bound to paper:1.5; no classification hypothesis is assumed.
\end{definition}

\begin{theorem}[The two supported square cases on the second line]
\label{sem:Main:two-line-known-squares}
\level{coarse}
\lean{KIP126.Main.adams_two_line_known_squares}
From \cite{LWX126}, Corollary 1.5.\par
For $j=5$ or $j=6$, the class $h_j^2\in E_2^{2,2^{j+1}}$ survives nontrivially to $E_\infty$ and belongs to the list $L$ above.
% Blueprint note: This is a locally completed proof from the existing upstream interfaces; their registered pending obligations remain. It proves neither the other listed survivors nor exhaustiveness. In particular it is not a proof binding for paper:1.5.
% Blueprint note: Detailed source attribution: (The $j=5,6$ cases of \cite{LWX126}, Corollary 1.5.)
\end{theorem}
\begin{proof}
For $j=5$, use the fixed Literature permanence input. For $j=6$, use the main permanence theorem. In either case the class is the product of the two standard $h_j$ classes, transported along the bidegree equation, and $j\leq6$.
\end{proof}

\begin{theorem}[The complete second-line classification in stem 126]
\label{sem:Main:two-line-stem126}
\level{coarse}
\lean{KIP126.Main.adams_two_line_stem126}
From \cite{LWX126}, Corollary 1.5.\par
For every $x\in E_2^{2,128}$, nonzero survival to $E_\infty$ is equivalent to $x=h_6^2$.
% Blueprint note: This is a complete classification in the single bidegree (2,128), conditional on the existing upstream obligations. It does not classify the entire second line.
% Blueprint note: Detailed source attribution: (The stem-126 specialization of \cite{LWX126}, Corollary 1.5.)
\end{theorem}
\begin{proof}
The fixed $E_2$ basis input says that $x=0$ or $x=h_6^2$. A zero $E_2$ image of a common permanent representative has zero $E_\infty$ image, so nonzero survival excludes the first case. Conversely, $h_6^2$ survives by the main theorem.
\end{proof}

\section{Choice transport and the inductive criterion}
\begin{remark}[Remark 7.4: normalization of the finite-quotient criterion]
\label{sem:Main:remark-7.4}
\source{AimPaper Remark 7.4}
\level{coarse}
From \cite{LWX126}, Remark 7.4.\par
Let $q_n$ denote the quotient map to $S^{0,0}/\lambda^n$. For every $a\in\pi_{125,130}S^{0,0}$ and $n\geq0$,
\[
 q_n(a)=0\quad\Longleftrightarrow\quad q_{n+1}(\lambda a)=0.
\]
In particular, with $a=\eta(\theta_5^{\mathrm{BX}})^2$ this identifies the two quotient conditions
\[
 q_r\bigl(\eta(\theta_5^{\mathrm{BX}})^2\bigr)=0
 \quad\Longleftrightarrow\quad
 q_{r+1}\bigl(\lambda\eta(\theta_5^{\mathrm{BX}})^2\bigr)=0.
\]
\lean{KIP126.Main.normalized_quotient}

% Blueprint note: This binds the public Main normalization theorem, whose proof calls the private finite-quotient argument. The quotient equivalence holds for every natural n; the cited fixed-choice BX criterion retains its positive-r hypothesis. No general vanishing of lambda torsion in other degrees is asserted.
\href{paper/main.tex.html#L2127}{Original remark}.
% Blueprint note: Detailed source attribution: (From \cite{LWX126}, Remark 7.4; the input criterion is \cite{BurklundXu}, Proposition 7.19.)
\end{remark}
\begin{proof}
Exactness of the lambda cofiber sequence identifies $\ker q_n$ with the image of $\lambda^n$. Multiplication by $\lambda$ is injective on $\pi_{125,130}S^{0,0}$ in the fixed system. Thus $\lambda a=\lambda^{n+1}b$ is equivalent to $a=\lambda^n b$. Specializing $a$ gives the change of quotient convention used in \ref{paper:7.3}.
\end{proof}

\begin{remark}[Remark 7.5: independence of the theta choice]
\label{sem:Main:remark-7.5}
\source{AimPaper Remark 7.5}
\level{coarse}
From \cite{LWX126}, Remark 7.5.\par
For every $\theta\in\pi_{62,64}S^{0,0}$ satisfying $\operatorname{ThetaChoice}(\theta)$,
\[
 \lambda\eta\theta^2=\lambda\eta(\theta_5^{\mathrm{BX}})^2.
\]
In particular, this expression has the same value for any two admissible choices. It is the expression, rather than $\theta$ itself, that is independent of the choice.
\lean{KIP126.Main.lambda_eta_theta_square_independent}

% Blueprint note: This binds the public Main choice-independence theorem. The filtration estimates, exponent-two identities and detector injectivity are actual upstream dependencies of that proof; neither the final permanence theorem nor an extra independence assumption is used.
\href{paper/main.tex.html#L2133}{Original remark}.
\end{remark}
\begin{proof}
The stem-62 choice comparison places $\theta-\theta_5^{\mathrm{BX}}$ in filtration at least six. The exponent-two square identity therefore puts $\theta^2-(\theta_5^{\mathrm{BX}})^2$ in filtration at least twelve. The established eta-error estimate raises the filtration of its eta multiple to at least fifteen. This eta-product difference is annihilated by two and has zero detector image, so the established filtration-fifteen detection theorem makes it zero. Multiplying by $\lambda$ gives the claimed equality.
\end{proof}

\begin{theorem}[The total boundary at every theta choice]
\label{sem:Main:boundary-formula}
\level{coarse}
\lean{KIP126.Main.delta_h6_square_eq_lambda_eta_theta_square}
From \cite{LWX126}, Remark 7.5.\par
For every admissible $\theta\in\pi_{62,64}S^{0,0}$, the total boundary of the first-quotient class $h_6^2$ satisfies
\[
 \delta_1(h_6^2)=\lambda\eta\theta^2.
\]
Here $\delta_1$ is the connecting map for the first lambda quotient. In particular, $\lambda\eta\theta^2$ is independent of the admissible theta choice.
\end{theorem}
\begin{proof}
Literature gives the boundary formula at $\theta_5^{\mathrm{BX}}$; \ref{sem:Main:remark-7.5} identifies its right-hand side with the expression at $\theta$.
\end{proof}

\begin{theorem}[Normalization and choice transport of the literature criterion]
\label{sem:Main:any-choice-criterion}
\level{coarse}
\lean{KIP126.Main.any_choice_criterion}
From \cite{LWX126}, Theorem 7.3 and Remarks 7.4--7.5; \cite{BurklundXu}, Proposition 7.19.\par
The literature criterion is recalled in \ref{paper:7.3}. On the fixed system its normalized form is as follows.
For every admissible $\theta$ and every natural number $r\geq0$,
\[
 h_6^2\text{ survives to }E_{r+3}
 \quad\Longleftrightarrow\quad
 [\lambda\eta\theta^2]=0\text{ in }S^{0,0}/\lambda^{r+1}.
\]
Moreover,
\[
 P\quad\Longleftrightarrow\quad\lambda\eta\theta^2=0
 \text{ in the untruncated synthetic sphere}.
\]
% Blueprint note: This is the fixed-system normalization and choice transport of an existing literature result. Its Lean proof calls the distinguished-choice Literature interfaces. The current final permanence proof calls those Literature interfaces directly and does not depend on this exported any-choice theorem.
\end{theorem}
\begin{proof}
For $r\geq1$, apply the distinguished-choice literature criterion and then combine \ref{sem:Main:remark-7.4} with \ref{sem:Main:remark-7.5}. At $r=0$, $h_6^2$ survives to $E_3$ and every lambda multiple maps to zero in $S^{0,0}/\lambda$. The untruncated equivalence is the distinguished-choice literature criterion followed by \ref{sem:Main:remark-7.5}.
\end{proof}

\begin{theorem}[Three exact forms for a theta-degree square]
\label{sem:Main:square-possibilities}
\level{coarse}
\uses{sem:Main:admissible-representatives,sem:Computation:p-q-correction,sem:Computation:correction-survival,sem:Computation:table-bases,sem:Literature:selected-lifts}
\lean{KIP126.Main.theta_square_three_possibilities}
From \cite{LWX126}, proof of Proposition 7.8.\par
For every $\theta\in\pi_{62,64}S^{0,0}$, at least one of the following holds:
\begin{enumerate}
\item $\theta^2=\lambda^6u$ for an admissible $u\in\pi_{124,134}S^{0,0}$;
\item $\theta^2=\lambda^9c$ for some $c\in\pi_{124,137}S^{0,0}$ nontrivially detected by the fixed correction class $e_0\Delta h_6g$ in bidegree $(13,137)$;
\item $\theta^2=\lambda^{10}a$ for some $a\in\pi_{124,138}S^{0,0}$.
\end{enumerate}
\end{theorem}

\section{Detection in finite lambda quotients}
\begin{definition}[The finite and untruncated detection predicates]
\label{sem:Main:finite-detection}
\level{coarse}
\lean{KIP126.Main.Paper.FiniteDetected,KIP126.Main.Paper.SphereDetected}
Let $x\in E_2^{s,t}$ and let $k\geq0$. For $n>0$ and
\[
 a\in\pi_{t-s,t-k}(S^{0,0}/\lambda^n),
\]
write $\operatorname{FiniteDetected}(n,s,t,k,x,a)$ if the quotient-projection page map sends the comparison of $\lambda^kx$ to a nonzero associated graded detector of $a$.
Write $\operatorname{SphereDetected}(s,t,k,x,a)$ for the corresponding nonzero detection in the untruncated sphere.
\end{definition}

\section{The retained contradiction argument}
\begin{remark}[The retained proof route for Proposition 7.9]
\label{sem:Main:weak-representative-route}
\level{coarse}
\lean{KIP126.Main.eta_extension_excluded_of_d6_vanishes}
\textit{(An alternative assembly of the argument in \cite{LWX126}, Proposition 7.9; compare Lemmas 7.11 and 7.16, Corollary 7.18, Proposition 7.20 and Fact 7.21.)}
Set $Q_n=S^{0,0}/\lambda^n$ and $Y=h_0^2x_{125,9,2}$.
Assume that $d_6(x_{126,8,4}+x_{126,8})=0$ and that there is an admissible
$u\in\pi_{124,134}S^{0,0}$ for which $\lambda^3\eta u$ is detected by
$\lambda^6h_1h_4x_{109,12}$. The conclusion remains exactly
\ref{paper:7.9}: these two conditions cannot hold together.

The original proof uses the detection of every member of a specified Toda
bracket in \ref{paper:7.16} and the exact $h_0$ relation for every detected
representative in \ref{paper:7.18}. The retained route instead starts with an
actual bracket member $z\in\pi_{125,132}Q_9$ and its exact target relation,
which yields $\lambda^2h_0z\ne0$. It removes the filtration-seven and
filtration-nine terms while preserving this multiplied value, obtaining
\[
 y\in F^{11}\pi_{125,132}Q_9,\qquad
 \lambda^2h_0y=\lambda^2h_0z\ne0.
\]
The finite-page calculations identify the detector of $y$ as $\lambda^4Y$.
There is no assertion that $y$ still belongs to the original Toda bracket,
or that every representative detected by $\lambda^4Y$ has the unmultiplied
exact extension of Corollary 7.18. The ninth-quotient $\nu$ extension from
\ref{paper:7.20}, the product calculations of \ref{paper:7.21}, and the
$\nu$-cofiber survival obstruction then give the contradiction.

The filtration-nine removal uses
\[
 \lambda^5\eta(\theta_5 b)=0\quad\text{in }\pi_{125,131}Q_9,
 \qquad b\text{ detected by }\Delta e_1+C_0+h_0^6h_5^2.
\]
This follows from $h_5^2b=0$, the stem-124 filtration-eleven and
filtration-twelve vanishings and the filtration-thirteen $h_1$ product
calculation already used in the proof of \ref{paper:7.11}, followed by the
finite-quotient filtration cutoff. It does not cancel $\eta$ from an
equality. The auxiliary $\langle\theta_5,2,b\rangle$ is constructed using
cofiber exactness, the quotient algebra, and the first-quotient product
comparison in place of the original Moss argument and synthetic lift check.

The relevant upstream nodes are \ref{sem:AdamsSynthetic:finite-quotient-separation},
\ref{sem:Literature:infinite-nu-window},
\ref{sem:Literature:infinite-finite-window},
\ref{sem:Literature:selected-lifts},
\ref{sem:Literature:quotient-algebras},
\ref{sem:Literature:first-quotient-products}, and
\ref{sem:Literature:may-tensor}. The finite-quotient separation statement has
a source proof route that does not assume the original spectrum is complete
(\cite{BHS}, Lemma A.14). The filtration--$\lambda$ comparison can be derived
from these finite windows, finite-quotient separation and cofiber exactness;
it is not a direct application of \cite{BHS}, Corollary A.19, to an ordinary
uncompleted sphere.
\uses{paper:7.11,paper:7.20,paper:7.21,sem:AdamsSynthetic:finite-quotient-separation,sem:Literature:infinite-nu-window,sem:Literature:infinite-finite-window,sem:Literature:infinite-lambda,sem:Literature:infinite-restriction,sem:Literature:selected-lifts,sem:Literature:quotient-algebras,sem:Literature:first-quotient-products,sem:Literature:may-tensor,sem:Literature:low-toda-comparison,sem:Literature:symmetric-toda-comparison,sem:Literature:toda-functoriality,sem:Literature:sphere-vanishing,sem:Computation:nondivisibility,sem:Computation:cofiber-results,sem:Computation:vanishing-regions,sem:Computation:product-boundaries,frontier:AdamsCalculations:e95db202825f,sem:ProofTools:sphere-label-detection,sem:ProofTools:finite-strip}
% Blueprint note: This entry states the same incompatibility conclusion as paper:7.9 and binds its existing public Main theorem. The remaining text explains its retained proof, without binding private helpers. The user retained this route after the upstream fidelity review. The two original conditional statements at paper:7.16 and paper:7.18 are now proved from this independently established incompatibility. Their direct local Toda and unmultiplied extension arguments are not asserted here. Main's local assembly still inherits the registered upstream obligations. A separate audit proved the complete fixed-model filtration--lambda comparison signature from existing upstream inputs, without using that target or Main, but its production proof remains pending. No public interface or fixed object was changed, and no mathematical novelty claim is made.
\end{remark}
